% GENERATED FILE. Edit canonical lessons, then run npm run generate:books.
% canonical-source-hash: fnv1a64:740d91feb92073d0
% canonical-lessons: FR-C422-ferry, FR-C422-traversee, FR-C422-cabine, FR-C422-gilet, FR-C422-debarquer

\chapter{La traversée}
\label{ch:fr-la-traversee}
\hlchaptermodality{422}

\begin{chapteropening}
\textbf{By the end of this chapter:} \emph{I can take a ferry: the ferry, the crossing, the cabin, the life jacket and getting off.}

Complete the last lesson of chapter 422: I can take a ferry: the ferry, the crossing, the cabin, the life jacket and getting off.
\end{chapteropening}

\section[le ferry]{le ferry — the ferry}

\label{lesson:FR-C422-ferry}



\begin{quote}
Before the new one: say the French for to cancel, then the French for the overbooking.
\end{quote}

\subsection*{You'll want to know: le ferry}
\textbf{le ferry} — \textquotedblleft{}the ferry\textquotedblright{}.

\begin{quote}
Nous prenons le ferry pour aller sur l'île.
\end{quote}

\begin{quote}
\emph{We take the ferry to go over to the island.}
\end{quote}

In Québec: \emph{le traversier}.

\begin{cousinweb}
\emph{ferry} was borrowed from English \emph{ferry} (EtymDB), from a verb meaning 'to carry over water'.
\end{cousinweb}

\begin{grammarlens}[title={What you've built}]
One more word for telling, explaining and describing at B1.
\end{grammarlens}

\subsection*{Your turn}
\begin{itemize}
\raggedright
  \item \emph{Say it:} \emph{le ferry}
  \item \emph{Say it:} \emph{le ferry}, once more
  \item \emph{Say it:} a sentence of your own with \emph{le ferry}
  \item \emph{Recall:} say the French for to cancel, then the French for the overbooking, then say \emph{le ferry} again
\end{itemize}

\begin{tcolorbox}[breakable,colback=teal!4,colframe=teal!35!black,title={Before you move on}]
What does le ferry mean? (\textquotedblleft{}The ferry\textquotedblright{}.) Say it once more.
\end{tcolorbox}
\section[la traversée]{la traversée — the crossing (by sea)}

\label{lesson:FR-C422-traversee}



\begin{quote}
Before the new one: say the French for the overbooking, then the French for the ferry.
\end{quote}

\subsection*{You'll want to know: la traversée}
\textbf{la traversée} — \textquotedblleft{}the crossing (by sea)\textquotedblright{}.

\begin{quote}
La traversée dure quatre heures.
\end{quote}

\begin{quote}
\emph{The crossing takes four hours.}
\end{quote}

Built on \emph{traverser} 'to cross'.

\begin{cousinweb}
\emph{traversée} is built on \emph{traverser} 'to cross', which goes back to Latin \emph{trānsversus} 'lying across', from \emph{vertere} 'to turn' (Lewis \& Short). EtymDB does not list the noun.
\end{cousinweb}

\begin{grammarlens}[title={What you've built}]
One more word for telling, explaining and describing at B1.
\end{grammarlens}

\subsection*{Your turn}
\begin{itemize}
\raggedright
  \item \emph{Say it:} \emph{la traversée}
  \item \emph{Say it:} \emph{la traversée}, once more
  \item \emph{Say it:} a sentence of your own with \emph{la traversée}
  \item \emph{Recall:} say the French for the overbooking, then the French for the ferry, then say \emph{la traversée} again
\end{itemize}

\begin{tcolorbox}[breakable,colback=teal!4,colframe=teal!35!black,title={Before you move on}]
What does la traversée mean? (\textquotedblleft{}The crossing (by sea)\textquotedblright{}.) Say it once more.
\end{tcolorbox}
\section[la cabine]{la cabine — the cabin}

\label{lesson:FR-C422-cabine}



\begin{quote}
Before the new one: say the French for the ferry, then the French for the crossing.
\end{quote}

\subsection*{You'll want to know: la cabine}
\textbf{la cabine} — \textquotedblleft{}the cabin\textquotedblright{}.

\begin{quote}
Pour la nuit, ils ont pris une cabine avec deux lits.
\end{quote}

\begin{quote}
\emph{For the night they took a cabin with two beds.}
\end{quote}

\emph{Une cabine d'essayage} — a fitting room.

\begin{cousinweb}
The origin of \emph{cabine} is not settled; it is usually linked to a northern dialect word and to English \emph{cabin}. EtymDB lists it without a source.
\end{cousinweb}

\begin{grammarlens}[title={What you've built}]
One more word for telling, explaining and describing at B1.
\end{grammarlens}

\subsection*{Your turn}
\begin{itemize}
\raggedright
  \item \emph{Say it:} \emph{la cabine}
  \item \emph{Say it:} \emph{la cabine}, once more
  \item \emph{Say it:} a sentence of your own with \emph{la cabine}
  \item \emph{Recall:} say the French for the ferry, then the French for the crossing, then say \emph{la cabine} again
\end{itemize}

\begin{tcolorbox}[breakable,colback=teal!4,colframe=teal!35!black,title={Before you move on}]
What does la cabine mean? (\textquotedblleft{}The cabin\textquotedblright{}.) Say it once more.
\end{tcolorbox}
\section[le gilet]{le gilet — the waistcoat; the cardigan; the life jacket (gilet de sauvetage)}

\label{lesson:FR-C422-gilet}



\begin{quote}
Before the new one: say the French for the crossing, then the French for the cabin.
\end{quote}

\subsection*{You'll want to know: le gilet}
\textbf{le gilet} — \textquotedblleft{}the waistcoat; the cardigan; the life jacket (gilet de sauvetage)\textquotedblright{}.

\begin{quote}
Sous chaque siège, il y a un gilet de sauvetage.
\end{quote}

\begin{quote}
\emph{Under every seat there is a life jacket.}
\end{quote}

\emph{Le gilet de sauvetage} — the life jacket. \emph{Un gilet} can also be a cardigan.

\begin{cousinweb}
\emph{gilet} goes back to Turkish \emph{yelek} 'waistcoat', which reached French through Arabic (EtymDB), probably by way of Spanish.
\end{cousinweb}

\begin{grammarlens}[title={What you've built}]
One more word for telling, explaining and describing at B1.
\end{grammarlens}

\subsection*{Your turn}
\begin{itemize}
\raggedright
  \item \emph{Say it:} \emph{le gilet}
  \item \emph{Say it:} \emph{le gilet}, once more
  \item \emph{Say it:} a sentence of your own with \emph{le gilet}
  \item \emph{Recall:} say the French for the crossing, then the French for the cabin, then say \emph{le gilet} again
\end{itemize}

\begin{tcolorbox}[breakable,colback=teal!4,colframe=teal!35!black,title={Before you move on}]
What does le gilet mean? (\textquotedblleft{}The waistcoat; the cardigan; the life jacket (gilet de sauvetage)\textquotedblright{}.) Say it once more.
\end{tcolorbox}
\section[débarquer]{débarquer — to get off (a boat or plane), to land}

\label{lesson:FR-C422-debarquer}



\begin{quote}
Before the new one: say the French for the cabin, then the French for the waistcoat.
\end{quote}

\subsection*{You'll want to know: débarquer}
\textbf{débarquer} — \textquotedblleft{}to get off (a boat or plane), to land\textquotedblright{}.

\begin{quote}
Les passagers débarquent sur le quai.
\end{quote}

\begin{quote}
\emph{The passengers get off onto the quay.}
\end{quote}

Informal: \emph{il débarque} — he turns up out of the blue.

\begin{cousinweb}
\emph{débarquer} is built from \emph{dé-} and \emph{barque} 'boat' (EtymDB), which goes back to Latin \emph{barca} (Lewis \& Short).
\end{cousinweb}

\begin{grammarlens}[title={What you've built}]
One more word for telling, explaining and describing at B1.
\end{grammarlens}

\subsection*{Your turn}
\begin{itemize}
\raggedright
  \item \emph{Say it:} \emph{débarquer}
  \item \emph{Say it:} \emph{débarquer}, once more
  \item \emph{Say it:} a sentence of your own with \emph{débarquer}
  \item \emph{Recall:} say the French for the cabin, then the French for the waistcoat, then say \emph{débarquer} again
\end{itemize}

\begin{tcolorbox}[breakable,colback=teal!4,colframe=teal!35!black,title={Before you move on}]
What does débarquer mean? (\textquotedblleft{}To get off (a boat or plane), to land\textquotedblright{}.) Say it once more.
\end{tcolorbox}
